\documentclass{article}
\usepackage{amsmath, amssymb, amsthm}
\usepackage{hyperref}
\usepackage{geometry}
\geometry{margin=1in}

\title{Erd\H{o}s Problem \#820}
\date{Last updated: 2025-08-31}
\author{Source: erdosproblems.com}

\begin{document}
\maketitle

\noindent\textbf{Status:} OPEN \quad \textbf{No prize}

\noindent\textbf{Tags:} number theory

\medskip

\noindent\textbf{Problem Statement:}

Let $H(n)$ be the smallest integer $l$ such that there exist $k&lt;l$ with $(k^n-1,l^n-1)=1$. Is it true that $H(n)=3$ infinitely often? (That is, $(2^n-1,3^n-1)=1$ infinitely often?) Estimate $H(n)$. Is it true that there exists some constant $c&gt;0$ such that, for all $\epsilon&#62;0$,\[H(n) &#62; \exp(n^{(c-\epsilon)/\log\log n})\]for infinitely many $n$ and\[H(n) &#60; \exp(n^{(c+\epsilon)/\log\log n})\]for all large enough $n$? Does a similar upper bound hold for the smallest $k$ such that $(k^n-1,2^n-1)=1$?




Erd\H{o}s \cite{Er74b} proved that there exists a constant $c&gt;0$ such that\[H(n) &gt; \exp(n^{c/(\log\log n)^2})\]for infinitely many $n$. van Doorn in the comments sketches a proof of the lower bound: that there exists some constant $c&gt;0$ and infinitely many $n$ such that\[H(n) &gt; \exp(n^{c/\log\log n}).\]The sequence $H(n)$ for $1\leq n\leq 10$ is\[3,3,3,6,3,18,3,6,3,12.\]The sequence of $n$ for which $(2^n-1,3^n-1)=1$ is A263647 in the OEIS. See also [770] .




References

[Er74b] Erd\H{o}s, P., Remarks on some problems in number theory . Math. Balkanica (1974), 197-202.

\noindent\textbf{Related OEIS sequences:} A263647


\bigskip
\noindent\small{Source: \url{https://www.erdosproblems.com/820}}
\end{document}
